% Generated draft: requires mathematical and source review.
\documentclass[11pt]{article}
\usepackage[margin=25mm]{geometry}
\usepackage{fontspec}
\setmainfont{Latin Modern Roman}
\usepackage{amsmath,amssymb,mathtools,unicode-math}
\setmathfont{Latin Modern Math}
\usepackage{enumitem,needspace,xurl,hyperref}
\setlength{\parindent}{0pt}
\setlength{\parskip}{5pt}
\setlength{\emergencystretch}{3em}
\newcommand{\N}{\mathbb{N}}
\newcommand{\Z}{\mathbb{Z}}
\newcommand{\Q}{\mathbb{Q}}
\newcommand{\R}{\mathbb{R}}
\newcommand{\C}{\mathbb{C}}
\newcommand{\F}{\mathbb{F}}
\newcommand{\A}{\mathbb{A}}
\newcommand{\PP}{\mathbb{P}}
\newcommand{\E}{\mathbb{E}}
\newcommand{\eps}{\varepsilon}
\newcommand{\dd}{\,\mathrm{d}}
\begin{document}
\section*{UnsolvedMath ID: 30003290}
\textbf{Source problem code:} catalogue entry
\subsection{Definitions and mathematical statement}
The supplied record asks for a maximal-monodromy statement for every integer \(d\geq 4\). It does not explicitly specify the ground field or the parameter space of the universal bundle. The formulation below spells out the customary interpretation using smooth complex plane curves and their projective parameter space; these conventions require confirmation against the cited report.

For an integer \(d\geq 4\), put
\[
r=d-3,\qquad g=\frac{(d-1)(d-2)}{2}.
\]
Under this interpretation, let
\[
U_d=\bigl\{[F]\in\PP(H^0(\PP^2_{\C},\mathcal O(d))): V(F)\text{ is smooth}\bigr\},
\]
and let
\[
\pi_d:\mathcal C_d\longrightarrow U_d,\qquad
\mathcal C_d=\{([F],x)\in U_d\times\PP^2_{\C}:F(x)=0\},
\]
be the universal family. Choose a base point \([F_0]\in U_d\), write \(C_0=V(F_0)\), and choose an orientation-preserving identification of the underlying surface of \(C_0\) with a closed connected oriented surface \(\Sigma_g\) of genus \(g\). Write
\[
\operatorname{Mod}_g=\pi_0\operatorname{Diff}^{+}(\Sigma_g).
\]
The bundle monodromy gives a homomorphism
\[
\rho_d:\pi_1(U_d,[F_0])\longrightarrow\operatorname{Mod}_g,
\qquad
\Gamma_d=\operatorname{im}(\rho_d).
\]

An \(r\)-spin structure here means an isomorphism class of a topological complex line bundle \(L\) together with an isomorphism \(L^{\otimes r}\cong K\), where \(K\) is the complex cotangent line bundle of the oriented surface equipped with a compatible complex structure. The associated topological classes carry the natural pullback action of \(\operatorname{Mod}_g\). Adjunction gives
\[
K_{C_0}\cong\mathcal O_{C_0}(d-3),
\]
so \(\mathcal O_{C_0}(1)\), with its induced \(r\)-th-power identification, determines the associated topological \(r\)-spin structure \(\phi_d\). Define its stabilizer by
\[
\operatorname{Mod}_g[\phi_d]
=\{[f]\in\operatorname{Mod}_g:f^*\phi_d=\phi_d\}.
\]
Both groups are expressed using the same identification with \(\Sigma_g\); changing that identification conjugates both groups simultaneously.

The proposed assertion is
\[
\text{for every integer }d\geq 4,\qquad
\Gamma_d=\operatorname{Mod}_g[\phi_d].
\]
The boundary case \(d=4\) is included, with \(r=1\). The complex parameter-space convention and the interpretation of \(\phi_d\) as a topological spin structure are explicit interpretive qualifications, rather than details verified from the report. The supplied record labels the assertion open; no proof or independent status determination is asserted here.
\subsection{Short English statement}
For every degree \(d\geq 4\), does monodromy of the universal family of smooth plane curves realize every orientation-preserving mapping class that preserves the associated \((d-3)\)-spin structure? The displayed formulation uses the customary complex parameter space and the spin structure induced by adjunction, subject to source confirmation.
\subsection{Sources}
\begin{itemize}
\item Supplied local UnsolvedMath registry record, permanent ID \(30003290\), title: Monodromy of Universal Plane-Curve Bundles; exact source problem code: \texttt{catalogue entry}. Collection: Oberwolfach Reports Open Problems. No individual proposer is supplied; the proposed year is recorded as 2016.
\item Supplied source citation: \emph{Surface Bundles} (2017), Oberwolfach Reports, pp. 3149--3195, DOI 10.4171/owr/2016/56. Supplied source URL: \url{https://doi.org/10.4171/owr/2016/56}. The background describes Oberwolfach Reports as documenting mathematical workshops and their open-problem discussions.
\item Supplied literature source: Oberwolfach Report 56 (2017), \url{https://doi.org/10.4171/owr/2016/56}. The literature assessment checked at \texttt{2026-08-21} labels the problem open, reports that the maximal-monodromy assertion is stated as a conjecture for all relevant degrees, and says that no complete proof was found. This is supplied dated curation derived from the OWR material, not independently verified here; time-sensitive status and cited claims require independent verification.
\item The registry assigns the statement status \texttt{corrected\_verified} and reports manual semantic, typographic, and LaTeX review of the clean statement. These are supplied editorial assessments, not an independent mathematical verification of this draft or of the conjecture.
\end{itemize}
\end{document}
